\pagestyle{empty}

%% レイアウト
\setlength{\topmargin}{-10.4mm}
\setlength{\headheight}{0mm}
\setlength{\headsep}{0mm}
\setlength{\textheight}{262mm}
\setlength{\textwidth}{180mm}
%\setlength{\topskip}{7mm}
\setlength{\evensidemargin}{-10.4mm} 
\setlength{\oddsidemargin}{-10.4mm} 
\setlength{\columnsep}{8mm}

%% パッケージ環境
% graphicx: 画像読み込み
%  dvipdfmxをドライバ指定することでpdf/png/jpeg形式の図を利用可能
%  異なるドライバを利用する場合はそのドライバ指定に変更
% \usepackage{graphicx} % suzuki
\usepackage[dvipdfmx]{graphicx} % suzuki
% subcaption: subfigureを置き換えたパッケージ（複数の図用）
% \setlength{\footskip}{12mm}
% \usepackage{subfigure}	 % suzuki
\usepackage{subcaption} % suzuki
% color：文章に色をつけたいとき
%  利用例：\textcolor{red}{文章}
% \usepackage{color}
\usepackage{booktabs}
\usepackage{setspace}

% 行間調整
\setstretch{0.9}

%sectionのフォントサイズ修正
\makeatletter
\def\section{\@startsection {section}{1}{\z@}{2.5ex plus -1ex minus -.2ex}{1.3 ex plus .1ex}{\large\bf}}
\makeatother 

%subsectionのフォントサイズ修正
\makeatletter
\def\subsection{\@startsection {subsection}{1}{\z@}{1.5ex plus -1ex minus -.4ex}{0.3 ex plus .1ex}{\bf}}
\makeatother 

% タイトル・著者設定コマンド
\makeatletter

% 情報を保存
\newcommand{\subtitle}[1]{\gdef\@subtitle{#1}}
\newcommand{\teacher}[1]{\gdef\@teacher{#1}}

% 初期値
\subtitle{}
\teacher{}

% タイトル生成
\renewcommand{\maketitle}{%
\twocolumn[
  \begin{center}
    {\LARGE\bfseries\@title\par}

    % サブタイトルが空でなければ表示
    \ifx\@subtitle\@empty\else
      \vspace{0.5em}
      {\Large\bfseries\@subtitle\par}
    \fi

    \vspace{1em}
    \@author\par
    \vspace{0.5em}
    \textbf{指導教員}\ \@teacher
  \end{center}
  \vspace{1em}
]
}

\makeatother
